% Apunts de teoria en PDF · 3r ESO · Unitat 4: Àlgebra. Polinomis
% Mateix estil que els apunts de Batxillerat (franja i caixes de teoria en blau, exemples en roig), però més visuals,
% curts i amb dibuixos, per a l'alumnat de 3r ESO. Es compila amb pdflatex dues vegades (per als enllaços de l'índex),
% amb logo-fentmates.pdf (_eines/logo/pdf/horitzontal-clar.pdf) al costat.
\documentclass[11pt,a4paper]{article}
\usepackage[utf8]{inputenc}
\usepackage[T1]{fontenc}
\usepackage[catalan]{babel}
\usepackage[sfdefault,book,lining]{FiraSans}
\usepackage[italic,noplusnominus,noequal]{mathastext}
\linespread{1.18}
\usepackage[a4paper,top=2.4cm,bottom=2.2cm,left=1.6cm,right=1.6cm,headheight=1.1cm,headsep=0.5cm,footskip=1cm]{geometry}
\usepackage{amsmath,amssymb,cancel}
\usepackage{icomma}
\usepackage{graphicx}
\usepackage[table]{xcolor}
\usepackage{tikz}\usetikzlibrary{calc,positioning,arrows.meta,decorations.pathreplacing,shapes.geometric}
\usepackage{fontawesome5}
\usepackage{ccicons}
\usepackage{fancyhdr}
\usepackage{enumitem}
\usepackage{array}
\usepackage{needspace}
\usepackage[most]{tcolorbox}
\usepackage{titlesec}
\usepackage[hidelinks]{hyperref}
\hypersetup{pdftitle={Apunts de teoria. Unitat 4: Àlgebra. Polinomis. 3r ESO},pdfsubject={Matemàtiques 3r ESO},pdfkeywords={llenguatge algebraic, valor numèric, monomis, polinomis, identitats notables, factor comú, divisió de polinomis, regla de Ruffini, teorema del residu}}

% colors de la web (bloc d'Àlgebra: blau)
\definecolor{fmnavy}{HTML}{1B2748}
\definecolor{fmblue}{HTML}{3D6AA0}
\definecolor{fmbluel}{HTML}{4E81BD}
\definecolor{fmyellow}{HTML}{FFC94D}
\definecolor{fmgold}{HTML}{E9A919}
\definecolor{fmaccent}{HTML}{FF7A50}
\definecolor{fmdef}{HTML}{1D4F73}\definecolor{fmdefbg}{HTML}{EEF4F9}
\definecolor{fmex}{HTML}{A3203A}\definecolor{fmexbg}{HTML}{FBF0F2}
\definecolor{fmsf}{HTML}{C2410C}
\definecolor{fmok}{HTML}{15803D}
\definecolor{fmmuted}{HTML}{5D6685}
\definecolor{fmsurface}{HTML}{EEF1F8}
\definecolor{fmA}{HTML}{2F6FD6}\definecolor{fmB}{HTML}{EA6A10}\definecolor{fmR}{HTML}{E11D48}\definecolor{fmK}{HTML}{16A34A}\definecolor{fmV}{HTML}{7C3AED}

% lletres
\newcommand{\comic}{\fontfamily{ComicNeue-TLF}\selectfont}
\newcommand{\heavy}{\fontseries{eb}\selectfont}
\newcommand{\semi}{\fontseries{sb}\selectfont}
\DeclareUnicodeCharacter{2212}{-}

\setlength{\parindent}{0pt}
\setlength{\parskip}{7pt}
\setlength{\lineskip}{4pt}\setlength{\lineskiplimit}{4pt}
\setlength{\emergencystretch}{3em}
\setlist{itemsep=3pt,topsep=3pt,leftmargin=1.4em}
\setcounter{secnumdepth}{1}
\renewcommand{\thesection}{4.\arabic{section}}
\titleformat{\section}[hang]{\color{fmnavy}\fontsize{23}{27}\heavy\selectfont}{\tcbox[on line,colback=fmblue,colframe=fmblue,arc=3mm,boxsep=2pt,left=6pt,right=6pt,top=3pt,bottom=3pt]{\color{white}\thesection}}{0.45em}{}[\vspace{5pt}{\color{fmblue}\titlerule[2pt]}]
\titlespacing*{\section}{0pt}{4pt}{6pt}
\titleformat{\subsection}{\color{fmblue}\Large\semi}{}{0pt}{\faCaretRight\ }
\titlespacing*{\subsection}{0pt}{8pt}{2pt}

% frase de cada apartat, «a mà»
\newcommand{\frase}[1]{\par\vspace{-2pt}{\comic\Large\color{fmbluel}\faCommentDots[regular]\ \ #1\par}\vspace{2pt}}
\newcommand{\seccio}[3]{\Needspace{16\baselineskip}\section{#2}\hypertarget{#1}{}\frase{#3}}
% notes a mà dins dels exemples
\newcommand{\nota}[1]{{\comic\small\color{fmex!85!black}#1}}
\newcommand{\notab}[1]{{\comic\small\color{fmblue}#1}}
% bé / malament
\newcommand{\be}[1]{{\color{fmok}\faCheckCircle}\ \ensuremath{#1}}
\newcommand{\mal}[1]{{\color{fmR}\faTimesCircle}\ {\color{fmR}\ensuremath{#1}}}
% pas numerat
\newcommand{\pas}[1]{\tikz[baseline=(p.base)]\node[circle,fill=fmblue,text=white,inner sep=1.2pt,minimum size=1.25em,font=\small\heavy](p){#1};}
% etiqueta de color
\newcommand{\chip}[2]{\tcbox[on line,colback=#1!14,colframe=#1!60,arc=2mm,boxsep=1pt,left=5pt,right=5pt,top=1pt,bottom=1pt]{#2}}
\newcommand{\fmkey}[1]{\colorbox{fmyellow!55}{\strut #1}}
\newcommand{\cA}[1]{{\color{fmA}#1}}\newcommand{\cB}[1]{{\color{fmB}#1}}\newcommand{\cK}[1]{{\color{fmK}#1}}\newcommand{\cR}[1]{{\color{fmR}#1}}\newcommand{\cV}[1]{{\color{fmV}#1}}

% caixes
\tcbset{fmbase/.style={enhanced,arc=3mm,boxrule=1.2pt,left=4mm,right=4mm,top=2mm,bottom=2.5mm,before skip=8pt,after skip=8pt,fonttitle=\heavy\large,before upper=\raggedright}}
\newtcolorbox{idea}[1][Idea clau]{fmbase,colback=fmdefbg,colframe=fmdef,coltitle=white,title={\faKey\ \ #1}}
\newtcolorbox{exemple}[1][Exemple]{fmbase,colback=fmexbg,colframe=fmex,colbacktitle=fmexbg,coltitle=fmex,title={\faPencil*\ \ #1},titlerule=0.5pt,titlerule style={fmex!40,dashed}}
\newtcolorbox{compte}[1][Compte!]{fmbase,colback=fmsf!6,colframe=fmsf,coltitle=white,fonttitle=\comic\bfseries\Large,title={\faExclamationTriangle\ \ #1}}
\newtcolorbox{truc}[1][Truc]{fmbase,colback=fmyellow!20,colframe=fmgold,colbacktitle=fmyellow,coltitle=fmnavy,fonttitle=\comic\bfseries\Large,title={\faLightbulb[regular]\ \ #1}}

% dibuix de la targeta de la unitat (com als apunts d'Àlgebra de Batxillerat), en blanc
\newcommand{\fmcard}{\begin{tikzpicture}[x=0.062cm,y=-0.062cm,white,line width=3pt,line cap=round]
  \begin{scope}[transparency group,opacity=0.15]
  \node[anchor=base west,inner sep=0pt,font=\fontsize{26}{26}\selectfont] at (14,44) {$P(x)=x^2+3x-4$};
  \draw[line width=2.4pt] (40,62) rectangle (84,106) (88,62) rectangle (106,106);
  \fill[opacity=0.3] (40,62) rectangle (84,106);
  \node[anchor=base west,inner sep=0pt,font=\fontsize{17}{17}\selectfont] at (52,90) {$x^2$};
  \node[anchor=base west,inner sep=0pt,font=\fontsize{17}{17}\selectfont] at (91,90) {$x$};
  \draw[line width=3.2pt] (140,106) .. controls (160,106) and (166,66) .. (186,66) .. controls (206,66) and (214,100) .. (226,74);
  \draw[line width=2pt,opacity=0.6] (134,106)--(230,106);
  \end{scope}
\end{tikzpicture}}
\newcommand{\fmbrand}{\raisebox{-0.25\height}{\includegraphics[height=0.5cm]{logo-fentmates.pdf}}}

% fitxes d'àlgebra (x², x, 1)
\newcommand{\tsq}[2]{\filldraw[fill=fmA!30,draw=fmA,line width=0.8pt,rounded corners=0.5pt] (#1,#2) rectangle ++(0.9,0.9); \node[font=\small] at ($(#1,#2)+(0.45,0.45)$) {$x^2$};}

% regla de Ruffini: #1 = a, #2 = coeficients, #3 = fila del mig (sense el primer), #4 = fila de baix, #5 = nombre de coeficients
\newcommand{\ruffini}[5]{\begin{tikzpicture}[x=1.25cm,y=0.85cm,every node/.style={font=\large}]
  \foreach \v [count=\i] in {#2} \node (t\i) at (\i,0) {$\v$};
  \foreach \v [count=\i] in {#3} {\pgfmathtruncatemacro\j{\i+1}\node (m\j) at (\j,-1) {$\v$};}
  \foreach \v [count=\i] in {#4} \node (b\i) at (\i,-2.1) {$\v$};
  \node[text=fmB] at (0,-1) {$#1$};
  \draw[line width=1.1pt,fmnavy] (0.45,0.5)--(0.45,-1.5) (0.45,-1.5)--(#5+0.5,-1.5);
  \node[draw=fmR,line width=1.3pt,rounded corners=2mm,minimum size=0.75cm] at (#5,-2.1) {};
  \draw[-{Latex[length=2.2mm]},fmK,line width=1.1pt] (t1.south) -- (b1.north);
  \pgfmathtruncatemacro\last{#5-1}
  \foreach \i in {1,...,\last} {\pgfmathtruncatemacro\j{\i+1}
    \draw[-{Latex[length=2.2mm]},fmB,line width=1.1pt,shorten >=1pt] (b\i.north east) -- (m\j.south west);
    \draw[-{Latex[length=2.2mm]},fmA!70,line width=1pt,dashed,shorten >=2pt] ([xshift=8pt]t\j.south) -- ([xshift=8pt]b\j.north);}
\end{tikzpicture}}

% capçalera i peu
\pagestyle{fancy}\fancyhf{}
\fancyhead[L]{\small\fmbrand}
\fancyhead[R]{\small\color{fmmuted}3r ESO · Unitat 4: Àlgebra. Polinomis}
\fancyfoot[L]{\scriptsize\color{fmmuted}\href{https://creativecommons.org/licenses/by-nc-sa/4.0/deed.ca}{\ccbyncsa}\ \ FentMates · fentmates.org · CC BY-NC-SA 4.0}
\fancyfoot[R]{\small\color{fmmuted}\thepage}
\renewcommand{\headrulewidth}{0.4pt}\renewcommand{\headrule}{\color{fmmuted!40}\hrule height \headrulewidth}
\fancypagestyle{primera}{\fancyhf{}\renewcommand{\headrulewidth}{0pt}
  \fancyfoot[L]{\scriptsize\color{fmmuted}\href{https://creativecommons.org/licenses/by-nc-sa/4.0/deed.ca}{\ccbyncsa}\ \ FentMates · fentmates.org · CC BY-NC-SA 4.0}
  \fancyfoot[R]{\small\color{fmmuted}\thepage}}

% targeta de l'índex
\newcommand{\itile}[4]{\begin{tcolorbox}[enhanced,colback=white,colframe=fmblue!35,arc=3mm,boxrule=1pt,boxsep=0pt,left=6pt,right=6pt,top=5pt,bottom=5pt,before skip=0pt,after skip=7pt]
  \hyperlink{#1}{\begin{minipage}[c]{0.9cm}\centering{\color{fmblue}\LARGE #3}\end{minipage}\hspace{6pt}%
  \begin{minipage}[c]{0.78\linewidth}{\color{fmblue}\heavy #2}\\[-2pt]{\color{fmnavy}\semi\large #4}\end{minipage}}\end{tcolorbox}}

\begin{document}
\thispagestyle{primera}
\begin{tikzpicture}[remember picture,overlay]
  \fill[fmblue] (current page.north west) rectangle ([yshift=-9.4cm]current page.north east);
  \begin{scope}
    \clip (current page.north west) rectangle ([yshift=-9.4cm]current page.north east);
    \node[anchor=north west,inner sep=0pt] at ([xshift=6.6cm,yshift=-1.3cm]current page.north west) {\fmcard};
  \end{scope}
  \node[anchor=north west,inner sep=0pt] at ([xshift=1.6cm,yshift=-1.1cm]current page.north west) {\tcbox[colback=white,colframe=white,arc=2.5mm,boxsep=0pt,left=7pt,right=7pt,top=5pt,bottom=5pt]{\includegraphics[height=0.75cm]{logo-fentmates.pdf}}};
\end{tikzpicture}
\vspace*{0.55cm}
\tcbox[colback=white!15!fmblue,colframe=white!15!fmblue,arc=2.5mm,boxsep=0pt,left=6pt,right=6pt,top=3pt,bottom=3pt]{\color{white}\small\bfseries Apunts de teoria · 3r ESO · Matemàtiques}

\vspace{2pt}
{\raggedright\color{white}\fontsize{30}{36}\heavy\selectfont{\color{fmyellow}Unitat 4}\\[2pt]Àlgebra. Polinomis\par}
\vspace{8pt}
{\color{white}\comic\Large Lletres que fan de nombres: com s'escriuen i com s'opera amb elles.\par}
\vspace*{\dimexpr 7.4cm-\pagetotal\relax}

{\color{fmnavy}\heavy\Large\faMapSigns\ \ Què aprendràs?}\par\smallskip
\begin{minipage}[t]{0.485\linewidth}
\itile{sec1}{4.1}{\faFont}{Llenguatge algebraic}
\itile{sec2}{4.2}{\faCube}{Monomis}
\itile{sec3}{4.3}{\faCalculator}{Operacions amb monomis}
\itile{sec4}{4.4}{\faLayerGroup}{Polinomis}
\itile{sec5}{4.5}{\faPlus}{Suma i resta de polinomis}
\end{minipage}\hfill
\begin{minipage}[t]{0.485\linewidth}
\itile{sec6}{4.6}{\faTimes}{Producte de polinomis}
\itile{sec7}{4.7}{\faStar}{Identitats notables}
\itile{sec8}{4.8}{\faUmbrella}{Extracció de factor comú}
\itile{sec9}{4.9}{\faDivide}{Divisió de polinomis}
\itile{sec10}{4.10}{\faBolt}{Regla de Ruffini}
\end{minipage}

\vspace{16pt}
{\color{fmnavy}\heavy\Large\faPalette\ \ Com es llegeixen aquests apunts}\par\smallskip
\begin{tcbraster}[raster columns=4,raster equal height,raster column skip=8pt,raster before skip=12pt]
\tcbox[colback=fmdefbg,colframe=fmdef,arc=3mm,boxrule=1.2pt,halign=center,width=\linewidth]{{\color{fmdef}\Large\faKey}\\[2pt]{\color{fmdef}\heavy Idea clau}\\[-1pt]{\small el que cal saber}}
\tcbox[colback=fmexbg,colframe=fmex,arc=3mm,boxrule=1.2pt,halign=center,width=\linewidth]{{\color{fmex}\Large\faPencil*}\\[2pt]{\color{fmex}\heavy Exemple}\\[-1pt]{\small pas a pas}}
\tcbox[colback=fmsf!6,colframe=fmsf,arc=3mm,boxrule=1.2pt,halign=center,width=\linewidth]{{\color{fmsf}\Large\faExclamationTriangle}\\[2pt]{\color{fmsf}\heavy Compte!}\\[-1pt]{\small errors típics}}
\tcbox[colback=fmyellow!20,colframe=fmgold,arc=3mm,boxrule=1.2pt,halign=center,width=\linewidth]{{\color{fmgold!80!black}\Large\faLightbulb[regular]}\\[2pt]{\color{fmnavy}\heavy Truc}\\[-1pt]{\small per a recordar-ho}}
\end{tcbraster}

% ============================================================ 4.1
\clearpage
\seccio{sec1}{Llenguatge algebraic}{Fem servir lletres per als nombres que no coneixem.}

\begin{minipage}[c]{0.27\linewidth}\centering
\begin{tikzpicture}
  \fill[fmyellow!75,draw=fmgold,line width=1.2pt,rounded corners=1.5mm] (0,0) rectangle (2.2,1.8);
  \fill[fmyellow!95!fmgold,draw=fmgold,line width=1.2pt,rounded corners=1mm] (-0.15,1.8) rectangle (2.35,2.25);
  \fill[fmaccent] (1.0,0) rectangle (1.2,2.25);
  \draw[fmaccent,line width=2pt] (1.1,2.25) .. controls (0.6,2.8) and (0.4,2.3) .. (1.1,2.25) .. controls (1.8,2.8) and (1.6,2.3) .. (1.1,2.25);
  \node[font=\fontsize{34}{34}\heavy\selectfont,text=fmnavy] at (0.52,0.9) {?};
  \node[font=\fontsize{34}{34}\heavy\selectfont,text=fmnavy] at (1.7,0.9) {?};
  \node[font=\fontsize{30}{30}\selectfont,text=fmblue] at (3.0,1.0) {$=x$};
\end{tikzpicture}
\end{minipage}\hfill
\begin{minipage}[c]{0.69\linewidth}
{\Large La lletra $x$ és com una \textbf{caixa} amb un \textbf{nombre amagat}.}\par
Pot valdre 2, 7, $-3$\dots{} Les lletres que fan de nombres s'anomenen \textbf{variables} (o \textbf{incògnites}, si volem trobar quant valen).
\begin{idea}
Una \textbf{expressió algebraica} barreja \cA{\textbf{nombres}}, \cB{\textbf{lletres}} i \cK{\textbf{operacions}}:\par\vspace{2pt}
{\centering\chip{fmblue}{$3x-5$}\ \ \chip{fmblue}{$2a+b^2$}\ \ \chip{fmblue}{$\dfrac{x+1}{2}$}\par}
\end{idea}
\end{minipage}

\subsection{De paraules a lletres}
\newcommand{\dic}[2]{#1 & {\large$#2$}\\}
% taula del diccionari dins d'una caixa amb les vores arrodonides (la capçalera és el títol de la caixa)
\newcommand{\dichead}{\begin{tabular}{@{}>{\raggedright}p{\dimexpr0.66\linewidth-12pt}>{\centering\arraybackslash}p{\dimexpr0.34\linewidth-12pt}@{}}\faComment[regular]\ Ho diem així\dots & \dots\ així\end{tabular}}
\newenvironment{dicbox}{\begin{tcolorbox}[enhanced,arc=3mm,boxrule=1.2pt,colframe=fmblue,colback=white,coltitle=white,fonttitle=\heavy,
  title={\dichead},left=0pt,right=0pt,top=0pt,bottom=1.5mm,boxsep=0pt,lefttitle=6pt,righttitle=6pt,before skip=0pt,after skip=0pt]\rowcolors{1}{white}{fmsurface}%
  \begin{tabular}{@{\hspace{6pt}}>{\raggedright}p{\dimexpr0.66\linewidth-12pt}>{\centering\arraybackslash}p{\dimexpr0.34\linewidth-12pt}@{\hspace{6pt}}}}
  {\end{tabular}\end{tcolorbox}}
\begin{minipage}[t]{0.485\linewidth}
\begin{dicbox}
\dic{El doble d'un nombre}{2x}
\dic{El triple d'un nombre}{3x}
\dic{La meitat d'un nombre}{\frac{x}{2}}
\dic{El quadrat d'un nombre}{x^2}
\dic{El cub d'un nombre}{x^3}
\dic{L'oposat d'un nombre}{-x}
\end{dicbox}
\end{minipage}\hfill
\begin{minipage}[t]{0.485\linewidth}
\begin{dicbox}
\dic{El següent d'un nombre}{x+1}
\dic{L'anterior d'un nombre}{x-1}
\dic{Un nombre parell}{2x}
\dic{Un nombre senar}{2x+1}
\dic{La suma de dos nombres}{x+y}
\dic{El producte de dos nombres}{xy}
\end{dicbox}
\end{minipage}

\begin{minipage}[t]{0.42\linewidth}
\begin{truc}
No escrivim el punt de multiplicar:\\ \mbox{$3\cdot x=3x$}. \ I l'1 tampoc: \mbox{$1x=x$}.
\end{truc}
\end{minipage}\hfill
\begin{minipage}[t]{0.55\linewidth}
\begin{compte}[Compte amb l'ordre!]
«El doble d'un nombre, més 3» \ $\to$\ \ \be{2x+3}\\[2pt]
«El doble de (un nombre més 3)» \ $\to$\ \ \be{2(x+3)}
\end{compte}
\end{minipage}

\subsection{Valor numèric}
{\large Canviem la lletra pel nombre i fem les operacions.} \ És com una \textbf{màquina}: hi entra un nombre i n'ix un altre.

\vspace{-10pt}
\begin{center}
\begin{tikzpicture}[scale=0.9,transform shape]
  \node[draw=fmblue,fill=fmdefbg,rounded corners=4mm,minimum width=4.6cm,minimum height=1.6cm,line width=1.4pt,font=\Large] (m) {$3x^2-5x+1$};
  \node[text=fmblue] at ([xshift=5mm,yshift=-4mm]m.north west) {\faCog};
  \node[text=fmblue] at ([xshift=-5mm,yshift=4mm]m.south east) {\faCog};
  \node[left=1.6cm of m,fill=fmyellow!60,rounded corners=2mm,font=\Large,inner sep=6pt] (in) {$x=-2$};
  \node[right=1.6cm of m,fill=fmK!20,draw=fmK,rounded corners=2mm,font=\Large\heavy,inner sep=6pt] (out) {$23$};
  \draw[-{Latex[length=3.5mm]},line width=1.8pt,fmmuted] (in) -- node[above,font=\comic\small]{entra} (m);
  \draw[-{Latex[length=3.5mm]},line width=1.8pt,fmmuted] (m) -- node[above,font=\comic\small]{ix} (out);
\end{tikzpicture}
\end{center}

\begin{minipage}[t]{0.53\linewidth}
\begin{exemple}[Exemple: \ $3x^2-5x+1$ \ per a \ $x=-2$]
{\large$3\cdot\cR{(-2)}^2-5\cdot\cR{(-2)}+1=12+10+1=\boxed{23}$}\\[2pt]
\nota{substituïm (amb parèntesis) i operem}
\end{exemple}
\end{minipage}\hfill
\begin{minipage}[t]{0.44\linewidth}
\begin{compte}
Els negatius, \textbf{sempre entre parèntesis}!\\[3pt]
\be{(-2)^2=4}\qquad \mal{-2^2=-4}\\[1pt]
{\small (són coses diferents)}
\end{compte}
\end{minipage}

% ============================================================ 4.2
\clearpage
\seccio{sec2}{Monomis}{Un monomi és un nombre multiplicat per lletres.}

\begin{center}
\begin{tikzpicture}
  \node[font=\fontsize{48}{48}\selectfont,inner sep=2pt] (c) {$\cA{-5}$};
  \node[font=\fontsize{48}{48}\selectfont,inner sep=2pt,right=2pt of c] (l) {$\cB{x^{\cK{2}}y^{\cK{1}}}$};
  \node[below left=0.9cm and -0.3cm of c,fill=fmA!14,draw=fmA,rounded corners=2mm,text=fmA,font=\heavy\large] (cl) {coeficient};
  \node[below right=0.9cm and -0.6cm of l,fill=fmB!14,draw=fmB,rounded corners=2mm,text=fmB,font=\heavy\large] (ll) {part literal};
  \node[right=2.4cm of l,fill=fmK!14,draw=fmK,rounded corners=2mm,text=fmK!80!black,font=\large,align=center] (g) {{\heavy grau}\\ $2+1=\mathbf{3}$};
  \draw[-{Latex[length=3mm]},fmA,line width=1.3pt] (cl) -- (c);
  \draw[-{Latex[length=3mm]},fmB,line width=1.3pt] (ll) -- (l);
  \draw[-{Latex[length=3mm]},fmK,line width=1.3pt] (g.west) to[bend right=25] ([yshift=-2pt]l.north east);
  \node[font=\comic,text=fmmuted,above=0pt of g] {suma els exponents};
  \node[font=\comic\small,text=fmmuted,below=2pt of g] {si no hi ha exponent, és 1: \ $y=y^1$};
\end{tikzpicture}
\end{center}

\begin{minipage}[c]{0.43\linewidth}
\begin{idea}
\begin{itemize}[leftmargin=1em]
\item \cA{\textbf{Coeficient}}: el nombre (amb el signe).
\item \cB{\textbf{Part literal}}: les lletres i els exponents.
\item \cK{\textbf{Grau}}: la suma dels exponents.
\end{itemize}
\end{idea}
\end{minipage}\hfill
\begin{minipage}[c]{0.54\linewidth}\centering\renewcommand{\arraystretch}{1.35}
% taula dins d'una caixa amb les vores arrodonides (la capçalera és el títol de la caixa)
\begin{tcolorbox}[enhanced,width=\dimexpr8cm+6\tabcolsep+2.4pt\relax,arc=3mm,boxrule=1.2pt,colframe=fmblue,colback=white,coltitle=white,fonttitle=\heavy,
  title={\begin{tabular}{@{}*{4}{w{c}{2cm}}@{}}Monomi & Coeficient & Part literal & Grau\end{tabular}},left=0pt,right=0pt,top=0pt,bottom=1mm,boxsep=0pt,lefttitle=0pt,righttitle=0pt,toptitle=0pt,bottomtitle=0pt,before skip=0pt,after skip=0pt]
\rowcolors{1}{white}{fmsurface}\begin{tabular}{@{}*{4}{w{c}{2cm}}@{}}
$7x^4$ & \cA{$7$} & \cB{$x^4$} & \cK{4}\\
$-a^2b^3$ & \cA{$-1$} & \cB{$a^2b^3$} & \cK{$2+3=5$}\\
$x$ & \cA{$1$} & \cB{$x$} & \cK{1}\\
$-6$ & \cA{$-6$} & \cB{---} & \cK{0}\\
\end{tabular}
\end{tcolorbox}
\end{minipage}

\begin{minipage}[t]{0.44\linewidth}
\begin{truc}
No veus el coeficient? És un \textbf{1}:\\
$x^2=\cA{1}x^2$ \qquad $-x=\cA{-1}x$\\[2pt]
Un nombre tot sol és un monomi de \textbf{grau 0}.
\end{truc}
\end{minipage}\hfill
\begin{minipage}[t]{0.53\linewidth}
\begin{compte}[Això NO són monomis]
\mal{\dfrac{3}{x}}\ \ {\small la lletra divideix}\\[3pt]
\mal{\sqrt{x}}\ \ {\small hi ha una arrel}\\[3pt]
\mal{2x+3}\ \ {\small és una suma (un polinomi)}
\end{compte}
\end{minipage}

\subsection{Monomis semblants: la mateixa família}
{\large Són \textbf{semblants} si tenen \textbf{la mateixa part literal}} (mateixes lletres i mateixos exponents). És com tindre el mateix cognom.

\begin{center}
\begin{tikzpicture}[card/.style={draw=#1,fill=#1!10,line width=1.2pt,rounded corners=3mm,minimum width=2.6cm,minimum height=1.25cm,font=\Large}]
  \node[card=fmK] (a) at (0,0) {$3\cB{x^2y}$};
  \node[card=fmK] (b) at (3.4,0) {$-\frac12\cB{x^2y}$};
  \node[font=\Large,text=fmK] at (1.7,0) {\faUsers};
  \node[text=fmK!80!black,font=\comic\large] at (1.7,-1.05) {\faCheckCircle\ sí que són semblants};
  \node[card=fmR] (c) at (8.4,0) {$3\cB{x^2y}$};
  \node[card=fmR] (d) at (11.8,0) {$3\cB{xy^2}$};
  \node[font=\Large,text=fmR] at (10.1,0) {$\neq$};
  \node[text=fmR,font=\comic\large] at (10.1,-1.05) {\faTimesCircle\ no ho són (exponents canviats)};
\end{tikzpicture}
\end{center}
\begin{idea}[Per a què serveix?]
\centering\Large Només es poden \textbf{sumar i restar} els monomis \textbf{semblants}.
\end{idea}

% ============================================================ 4.3
\clearpage
\seccio{sec3}{Operacions amb monomis}{Per a sumar, han de ser de la mateixa família. Per a multiplicar, no cal.}

\subsection{Suma i resta}
\begin{minipage}[c]{0.48\linewidth}
\begin{tikzpicture}[font=\Large]
  \node at (0,0) {3 \textcolor{fmR}{\faIcon{apple-alt}}};
  \node at (1.15,0) {$+$};
  \node at (2.3,0) {2 \textcolor{fmR}{\faIcon{apple-alt}}};
  \node at (3.45,0) {$=$};
  \node at (4.6,0) {\textbf{5} \textcolor{fmR}{\faIcon{apple-alt}}};
  \node[text=fmok] at (5.7,0) {\faCheckCircle};
  \node at (0,-1.1) {3 \textcolor{fmR}{\faIcon{apple-alt}}};
  \node at (1.15,-1.1) {$+$};
  \node at (2.3,-1.1) {2 \textcolor{fmK}{\faLemon}};
  \node at (3.45,-1.1) {$\to$};
  \node[font=\comic\normalsize,text=fmR,align=center] at (4.75,-1.1) {no es poden\\ajuntar};
  \node[text=fmR] at (5.95,-1.1) {\faTimesCircle};
\end{tikzpicture}
\end{minipage}\hfill
\begin{minipage}[c]{0.5\linewidth}
\begin{tikzpicture}[x=1cm,y=1cm]
  \foreach \i in {0,...,4} {\tsq{\i*1}{0}}
  \node[font=\Large] at (5.25,0.45) {$+$};
  \foreach \i in {0,...,2} {\tsq{5.6+\i*1}{0}}
  \node[font=\Large,anchor=west] at (0,-0.65) {$5x^2+3x^2=\mathbf{8x^2}$};
  \node[font=\comic,text=fmmuted,anchor=west] at (4.6,-0.65) {8 peces iguals};
\end{tikzpicture}
\end{minipage}

\begin{minipage}[t]{0.52\linewidth}
\begin{idea}
\textbf{Sumem els coeficients}. La part literal es queda \textbf{igual}.\\[4pt]
$5x^2+3x^2=(5+3)x^2=8x^2$\\[2pt]
$4ab-7ab=-3ab$\\[2pt]
$5x^2+3x$ \ \notab{no són semblants: es queda així}
\end{idea}
\end{minipage}\hfill
\begin{minipage}[t]{0.45\linewidth}
\begin{compte}[Error típic]
\mal{5x^2+3x^2=8x^4}\\[3pt]
\be{5x^2+3x^2=8x^2}\\[3pt]
Quan \textbf{sumem}, els exponents \textbf{no canvien}.
\end{compte}
\end{minipage}

\subsection{Producte, quocient i potència}
\begin{tcbraster}[raster columns=3,raster equal height,raster column skip=8pt]
\begin{tcolorbox}[fmbase,colback=fmdefbg,colframe=fmdef,coltitle=white,title={\faTimes\ \ Producte},before upper=\centering]
Nombres: \textbf{multipliquem}\\ Exponents: \fmkey{\textbf{sumem}}\\[3pt]
{\large$x^m\cdot x^n=x^{m+n}$}\\[4pt]
$(-3x^2y)\cdot(4xy^3)=-12x^3y^4$
\end{tcolorbox}
\begin{tcolorbox}[fmbase,colback=fmdefbg,colframe=fmdef,coltitle=white,title={\faDivide\ \ Quocient},before upper=\centering]
Nombres: \textbf{dividim}\\ Exponents: \fmkey{\textbf{restem}}\\[3pt]
{\large$x^m:x^n=x^{m-n}$}\\[4pt]
$15x^5y^2:(3x^2y)=5x^3y$
\end{tcolorbox}
\begin{tcolorbox}[fmbase,colback=fmdefbg,colframe=fmdef,coltitle=white,title={\faSuperscript\ \ Potència},before upper=\centering]
Nombre: \textbf{l'elevem}\\ Exponents: \fmkey{\textbf{multipliquem}}\\[3pt]
{\large$(x^m)^n=x^{m\cdot n}$}\\[4pt]
$(2x^3)^2=4x^6$
\end{tcolorbox}
\end{tcbraster}

\begin{minipage}[t]{0.62\linewidth}
\begin{exemple}
\vspace{-6pt}
\begin{align*}
2x^3\cdot5x^4&=(\cA{2\cdot5})\,x^{\cK{3+4}}=\boxed{10x^7}\\
20x^6y^3:(-4x^2y)&=(\cA{20:(-4)})\,x^{\cK{6-2}}y^{\cK{3-1}}=\boxed{-5x^4y^2}
\end{align*}
\nota{\cA{nombres amb nombres} \quad\cK{exponents de cada lletra}}
\end{exemple}
\end{minipage}\hfill
\begin{minipage}[t]{0.35\linewidth}
\begin{truc}
Recorda: \ $x^0=1$\\[4pt]
$6x^2y:(2x^2)=3\,\cR{x^0}\,y$\\[2pt]
$\phantom{6x^2y:(2x^2)}=3\cdot\cR{1}\cdot y=3y$
\end{truc}
\end{minipage}

% ============================================================ 4.4
\clearpage
\seccio{sec4}{Polinomis}{Un polinomi és una suma de monomis (no semblants).}

\begin{center}
\chip{fmblue}{\semi Monomi} \ 1 terme: $5x$ \qquad
\chip{fmblue}{\semi Binomi} \ 2 termes: $2x+3$ \qquad
\chip{fmblue}{\semi Trinomi} \ 3 termes: $x^2-5x+6$
\end{center}

\begin{center}
\begin{tikzpicture}[t/.style={font=\fontsize{30}{30}\selectfont,inner sep=3pt,text height=25pt,text depth=8pt}]
  \node[t] (P) {$P(x)=$};
  \node[t,right=0pt of P,fill=fmA!12,rounded corners=2mm] (a) {$\cA{4}x^{\cK{3}}$};
  \node[t,right=5pt of a] (b) {$-2x^2$};
  \node[t,right=2pt of b] (c) {$+x$};
  \node[t,right=2pt of c,fill=fmB!14,rounded corners=2mm] (d) {$\cB{-7}$};
  \node[above left=0.7cm and -0.6cm of a,fill=fmA!14,draw=fmA,rounded corners=2mm,font=\large,align=center] (cp) {{\heavy\color{fmA}coeficient principal}: 4};
  \coordinate (e3) at ([xshift=-10pt,yshift=-19pt]a.north east);% sota l'exponent 3
  \node[anchor=north,fill=fmK!14,draw=fmK,rounded corners=2mm,font=\large,align=center] (gr) at ([yshift=-0.8cm]e3 |- a.south) {{\heavy\color{fmK!80!black}grau del polinomi}: 3\\[-2pt]{\comic\small el grau més alt}};
  \node[above right=0.7cm and -0.8cm of d,fill=fmB!14,draw=fmB,rounded corners=2mm,font=\large,align=center] (ti) {{\heavy\color{fmB}terme independent}: $-7$\\[-2pt]{\comic\small el que no té lletra}};
  \draw[-{Latex[length=3mm]},fmA,line width=1.3pt] (cp.south -| {$(a.west)+(11pt,0)$}) -- ([xshift=11pt,yshift=-1pt]a.north west);
  \draw[-{Latex[length=3mm]},fmK,line width=1.3pt] (gr.north) -- (e3);
  \draw[-{Latex[length=3mm]},fmB,line width=1.3pt] (ti) -- (d.north);
\end{tikzpicture}
\end{center}

\begin{minipage}[t]{0.38\linewidth}
\begin{idea}[Ordenat i complet]
\textbf{Ordenat}: de major a menor grau.\\[3pt]
\textbf{Complet}: no en falta cap grau (fins al terme independent).
\end{idea}
\end{minipage}\hfill
\begin{minipage}[t]{0.59\linewidth}
\begin{exemple}[Exemple: \ $x^4-3x+1$]
\centering
\begin{tikzpicture}[s/.style={font=\Large,minimum height=0.9cm,minimum width=1.3cm,rounded corners=2mm}]
  \node[s,draw=fmblue,fill=fmdefbg] at (0,0) {$x^4$};
  \node[s,draw=fmR,dashed,line width=1pt,text=fmR] at (1.55,0) {?\,$x^3$};
  \node[s,draw=fmR,dashed,line width=1pt,text=fmR] at (3.1,0) {?\,$x^2$};
  \node[s,draw=fmblue,fill=fmdefbg] at (4.65,0) {$-3x$};
  \node[s,draw=fmblue,fill=fmdefbg] at (6.2,0) {$+1$};
\end{tikzpicture}\\[3pt]
\raggedright Està ordenat \be{}, però \textbf{no és complet}: falten $x^3$ i $x^2$. \ \nota{Ho podem escriure: $x^4+0x^3+0x^2-3x+1$}
\end{exemple}
\end{minipage}

\subsection{Valor numèric i arrels}
{\large $P(a)$ vol dir: canvia la $x$ per $a$ i calcula.}

\begin{minipage}[c]{0.45\linewidth}
\begin{idea}
Si \ $P(a)=0$, \ diem que $a$ és una \fmkey{\textbf{arrel}} del polinomi.
\end{idea}
\end{minipage}\hfill
\begin{minipage}[c]{0.52\linewidth}
\begin{exemple}[Exemple: \ $P(x)=x^2-5x+6$]
$P(\cR{1})=1-5+6=2$\\[3pt]
$P(\cR{2})=4-10+6=\boxed{0}$ \ \ {\color{fmok}\Large\faBullseye}\ \ \nota{2 és una arrel!}
\end{exemple}
\end{minipage}

% ============================================================ 4.5
\clearpage
\seccio{sec5}{Suma i resta de polinomis}{Ajuntem els termes de la mateixa família.}

\begin{idea}
\begin{tabular}{@{}p{0.08\linewidth}p{0.84\linewidth}@{}}
{\color{fmblue}\Huge\faPlus} & \textbf{Suma}: posem un polinomi al costat de l'altre i \textbf{reduïm els semblants}.\\[6pt]
{\color{fmblue}\Huge\faMinus} & \textbf{Resta}: \textbf{canviem el signe de TOTS} els termes del segon i sumem.
\end{tabular}
\end{idea}

{\large Ho fem en \textbf{columnes}: cada columna, una família.} \quad $P(x)=3x^3-2x^2+5$ \quad $Q(x)=x^3+4x^2-3x-1$

\begin{center}\renewcommand{\arraystretch}{1.5}\large
\begin{tabular}{r>{\columncolor{fmV!10}}r>{\columncolor{fmA!10}}r>{\columncolor{fmK!10}}r>{\columncolor{fmB!10}}r}
 & {\comic\normalsize\color{fmV} els $x^3$} & {\comic\normalsize\color{fmA} els $x^2$} & {\comic\normalsize\color{fmK} les $x$} & {\comic\normalsize\color{fmB} els nombres}\\
$P(x)$: & $3x^3$ & $-2x^2$ & & $+5$\\
$Q(x)$: & $x^3$ & $+4x^2$ & $-3x$ & $-1$\\ \hline
$P+Q$: & $\mathbf{4x^3}$ & $\mathbf{+2x^2}$ & $\mathbf{-3x}$ & $\mathbf{+4}$
\end{tabular}
\end{center}

\subsection{Per a restar: el signe menys «entra» al parèntesi}
\begin{center}
\begin{tikzpicture}[t/.style={font=\fontsize{22}{22}\selectfont,inner sep=2pt,text height=18pt,text depth=6pt}]
  \node[t,text=fmR] (m) {$-$};
  \node[t,right=0pt of m] (p1) {$\big($};
  \node[t,right=0pt of p1,minimum width=1.65cm] (t1) {$x^3$};
  \node[t,right=0pt of t1,minimum width=1.65cm] (t2) {$+4x^2$};
  \node[t,right=0pt of t2,minimum width=1.65cm] (t3) {$-3x$};
  \node[t,right=0pt of t3,minimum width=1.65cm] (t4) {$-1$};
  \node[t,right=0pt of t4] (p2) {$\big)$};
  \foreach \k/\b in {t1/50,t2/45,t3/40,t4/35} \draw[-{Latex[length=2.5mm]},fmR,line width=1.1pt] (m.north) to[bend left=\b] (\k.north);
  \node[t,below=1.0cm of t1,text=fmR] (r1) {$-x^3$};
  \node[t,below=1.0cm of t2,text=fmR] (r2) {$-4x^2$};
  \node[t,below=1.0cm of t3,text=fmR] (r3) {$+3x$};
  \node[t,below=1.0cm of t4,text=fmR] (r4) {$+1$};
  \foreach \i in {1,...,4} \draw[-{Latex[length=2.5mm]},fmmuted,line width=1pt] (t\i) -- (r\i);
  \node[t] at (m |- r1) {$=$};
  \node[right=1.2cm of p2,font=\comic\large,text=fmR,align=left] {canvien\\TOTS els signes};
\end{tikzpicture}
\end{center}

\begin{minipage}[t]{0.58\linewidth}
\begin{exemple}[Exemple: \ $P-Q$]
\vspace{-6pt}
\begin{align*}
P-Q&=3x^3-2x^2+5\,\cR{-\,x^3-4x^2+3x+1} && \text{\nota{canviem signes}}\\
&=\boxed{2x^3-6x^2+3x+6} && \text{\nota{reduïm}}
\end{align*}
\end{exemple}
\end{minipage}\hfill
\begin{minipage}[t]{0.39\linewidth}
\begin{compte}[L'error més freqüent]
Canviar només el \textbf{primer} signe:\\[3pt]
\mal{-(x^3+4x^2\dots)=-x^3+4x^2\dots}\\[3pt]
El menys afecta \textbf{tot} el parèntesi.
\end{compte}
\end{minipage}

% ============================================================ 4.6
\clearpage
\seccio{sec6}{Producte de polinomis}{Cada terme multiplica TOTS els termes de l'altre.}

\subsection{Monomi per polinomi}
\begin{center}
\begin{tikzpicture}[t/.style={font=\fontsize{22}{22}\selectfont,inner sep=2pt,text height=18pt,text depth=6pt}]
  \node[t,text=fmB] (m) {$-2x$};
  \node[t,right=0pt of m] (p1) {$\big($};
  \node[t,right=0pt of p1] (t1) {$3x^2$};
  \node[t,right=4pt of t1] (t2) {$-x$};
  \node[t,right=4pt of t2] (t3) {$+4$};
  \node[t,right=0pt of t3] (p2) {$\big)$};
  \foreach \k/\b/\c in {t1/55/fmV,t2/50/fmA,t3/45/fmK} \draw[-{Latex[length=2.5mm]},\c,line width=1.3pt] (m.north) to[bend left=\b] (\k.north);
  \node[t,right=0.3cm of p2] (eq) {$=$};
  \node[t,right=0.3cm of eq] {$\cV{-6x^3}\cA{+2x^2}\cK{-8x}$};
\end{tikzpicture}
\end{center}

\subsection{Polinomi per polinomi: amb una taula}
\begin{minipage}[c]{0.5\linewidth}\centering
\begin{tikzpicture}[x=1.75cm,y=0.95cm,c/.style={font=\large}]
  \begin{scope}
  \clip[rounded corners=3mm] (0,0) rectangle (4,3);
  \fill[fmdefbg] (0,0) rectangle (1,3) (1,2) rectangle (4,3);
  \fill[fmV!18] (1,1) rectangle (2,2);
  \fill[fmA!18] (2,1) rectangle (3,2);
  \fill[fmK!18] (3,1) rectangle (4,2);
  \fill[fmA!18] (1,0) rectangle (2,1);
  \fill[fmK!18] (2,0) rectangle (3,1);
  \fill[fmB!18] (3,0) rectangle (4,1);
  \end{scope}
  \foreach \x in {1,2,3} \draw[fmblue,line width=0.8pt] (\x,0)--(\x,3);
  \foreach \y in {1,2} \draw[fmblue,line width=0.8pt] (0,\y)--(4,\y);
  \draw[fmblue,line width=1.2pt,rounded corners=3mm] (0,0) rectangle (4,3);
  \node[c,font=\large\heavy,text=fmblue] at (0.5,2.5) {$\times$};
  \foreach \h [count=\i] in {x^2,+4x,-1} \node[c,font=\large\heavy,text=fmB] at (\i+0.5,2.5) {$\h$};
  \foreach \h [count=\j] in {2x,-3} \node[c,font=\large\heavy,text=fmB] at (0.5,2.5-\j) {$\h$};
  \node[c] at (1.5,1.5) {$2x^3$};  \node[c] at (2.5,1.5) {$8x^2$};   \node[c] at (3.5,1.5) {$-2x$};
  \node[c] at (1.5,0.5) {$-3x^2$}; \node[c] at (2.5,0.5) {$-12x$};  \node[c] at (3.5,0.5) {$+3$};
\end{tikzpicture}\\[2pt]
{\comic\color{fmmuted}$(2x-3)\cdot(x^2+4x-1)$}
\end{minipage}\hfill
\begin{minipage}[c]{0.47\linewidth}
\pas{1} Cada casella: \textbf{fila $\times$ columna}.\\[4pt]
\pas{2} Ajuntem les caselles del \textbf{mateix color} (són semblants):\\[4pt]
\hspace*{1.5em}$\cV{2x^3}\ \cA{+8x^2-3x^2}\ \cK{-2x-12x}\ \cB{+3}$\\[6pt]
\pas{3} Resultat:\\[2pt]
\hspace*{1.5em}{\Large$\boxed{2x^3+5x^2-14x+3}$}
\end{minipage}

\medskip
\begin{minipage}[t]{0.5\linewidth}
\begin{exemple}[Exemple: \ $(x+4)(x-6)$]
\vspace{-6pt}
\begin{align*}
(x+4)(x-6)&=x^2-6x+4x-24 && \text{\nota{tots per tots}}\\
&=\boxed{x^2-2x-24} && \text{\nota{reduïm}}
\end{align*}
\end{exemple}
\end{minipage}\hfill
\begin{minipage}[t]{0.47\linewidth}
\begin{truc}
El \textbf{grau del producte} és la \textbf{suma dels graus}:\\[3pt]
grau 1 \ $\times$ \ grau 2 \ $\to$ \ grau \textbf{3}
\end{truc}
\end{minipage}

\subsection{Polinomi per polinomi: terme a terme}
\begin{center}
\begin{tikzpicture}[t/.style={font=\fontsize{22}{22}\selectfont,inner sep=2pt,text height=18pt,text depth=6pt}]
  \node[t] (q1) {$\big($};
  \node[t,right=0pt of q1,text=fmA] (f1) {$2x$};
  \node[t,right=4pt of f1,text=fmB] (f2) {$+3$};
  \node[t,right=0pt of f2] (q2) {$\big)$};
  \node[t,right=0pt of q2] (p1) {$\big($};
  \node[t,right=0pt of p1] (t1) {$x^2$};
  \node[t,right=4pt of t1] (t2) {$-x$};
  \node[t,right=4pt of t2] (t3) {$+4$};
  \node[t,right=0pt of t3] (p2) {$\big)$};
  \foreach \k/\b in {t1/50,t2/45,t3/40} \draw[-{Latex[length=2.5mm]},fmA,line width=1.3pt] (f1.north) to[bend left=\b] (\k.north);
  \foreach \k/\b in {t1/50,t2/45,t3/40} \draw[-{Latex[length=2.5mm]},fmB,line width=1.3pt] (f2.south) to[bend right=\b] (\k.south);
  \node[font=\comic,text=fmA] at ([yshift=1.35cm]t2.north) {el $2x$ per tots};
  \node[font=\comic,text=fmB] at ([yshift=-1.35cm]t2.south) {el $+3$ per tots};
  \node[anchor=west,font=\Large] (r1) at ([xshift=0.5cm,yshift=0.5cm]p2.east) {$=\cA{2x^3-2x^2+8x}\,\cB{+3x^2-3x+12}$};
  \node[anchor=north west,font=\Large,inner sep=0pt] (r2) at ([yshift=-8pt]r1.south west) {\hspace{3.333pt}$=\boxed{2x^3+x^2+5x+12}$};
  \node[anchor=west,font=\comic\small,text=fmex!85!black] at (r2.east) {\ \ reduïm els semblants};
\end{tikzpicture}
\end{center}

% ============================================================ 4.7
\clearpage
\seccio{sec7}{Identitats notables}{Tres multiplicacions famoses: aprèn-te-les!}

\begin{tcbraster}[raster columns=3,raster equal height,raster column skip=8pt]
\begin{tcolorbox}[fmbase,colback=fmdefbg,colframe=fmdef,coltitle=white,title={\faStar\ \ Quadrat d'una suma},before upper=\centering,fonttitle=\heavy\small,left=2mm,right=2mm]
\vspace{4pt}{\large$(\cA{a}+\cB{b})^2=\cA{a}^2\,{+}\,2\cA{a}\cB{b}+\cB{b}^2$}
\end{tcolorbox}
\begin{tcolorbox}[fmbase,colback=fmdefbg,colframe=fmdef,coltitle=white,title={\faStar\ \ Quadrat d'una diferència},before upper=\centering,fonttitle=\heavy\small,left=2mm,right=2mm]
\vspace{4pt}{\large$(\cA{a}-\cB{b})^2=\cA{a}^2\,{-}\,2\cA{a}\cB{b}+\cB{b}^2$}
\end{tcolorbox}
\begin{tcolorbox}[fmbase,colback=fmdefbg,colframe=fmdef,coltitle=white,title={\faStar\ \ Suma per diferència},before upper=\centering,fonttitle=\heavy\small,left=2mm,right=2mm]
\vspace{4pt}{\large$(\cA{a}+\cB{b})(\cA{a}-\cB{b})=\cA{a}^2-\cB{b}^2$}
\end{tcolorbox}
\end{tcbraster}
\begin{center}
{\comic\large\color{fmblue}«El quadrat del primer, més el doble del primer pel segon, més el quadrat del segon»}
\end{center}

\subsection{Ho veiem amb dibuixos}
\begin{minipage}[c]{0.36\linewidth}\centering
\begin{tikzpicture}[x=1cm,y=1cm]
  \def\a{2.5}\def\b{1.1}
  \fill[fmA!28] (0,\b) rectangle (\a,\a+\b); \node[font=\Large] at (\a/2,\b+\a/2) {$\cA{a}^2$};
  \fill[fmK!25] (\a,\b) rectangle (\a+\b,\a+\b); \node at (\a+\b/2,\b+\a/2) {$\cA{a}\cB{b}$};
  \fill[fmK!25] (0,0) rectangle (\a,\b); \node at (\a/2,\b/2) {$\cA{a}\cB{b}$};
  \fill[fmB!28] (\a,0) rectangle (\a+\b,\b); \node at (\a+\b/2,\b/2) {$\cB{b}^2$};
  \draw[fmnavy,line width=1pt] (0,0) rectangle (\a+\b,\a+\b) (\a,0)--(\a,\a+\b) (0,\b)--(\a+\b,\b);
  \draw[decorate,decoration={brace,amplitude=5pt},fmA,thick] (0,\a+\b+0.08)--(\a,\a+\b+0.08) node[midway,above=5pt,text=fmA]{$a$};
  \draw[decorate,decoration={brace,amplitude=5pt},fmB,thick] (\a,\a+\b+0.08)--(\a+\b,\a+\b+0.08) node[midway,above=5pt,text=fmB]{$b$};
\end{tikzpicture}\\[3pt]
{\comic\color{fmmuted}$(a+b)^2=a^2+\mathbf{2}ab+b^2$}
\end{minipage}\hfill
\begin{minipage}[c]{0.62\linewidth}\centering
\begin{tikzpicture}[x=1cm,y=1cm]
  \def\a{2.5}\def\b{1.0}
  \fill[fmA!28] (0,0) rectangle (\a,\a-\b);
  \fill[fmK!25] (0,\a-\b) rectangle (\a-\b,\a);
  \draw[fmB,dashed,line width=1pt] (\a-\b,\a-\b) rectangle (\a,\a); \node[text=fmB] at (\a-\b/2,\a-\b/2) {$b^2$};
  \draw[fmnavy,line width=1pt] (0,0)--(\a,0)--(\a,\a-\b)--(\a-\b,\a-\b)--(\a-\b,\a)--(0,\a)--cycle;
  \node[font=\comic,text=fmmuted] at (\a/2,-0.4) {$a^2-b^2$};
  \draw[-{Latex[length=3mm]},fmmuted,line width=1.4pt] (\a+0.4,\a/2) -- node[above,font=\comic\small,text=fmmuted,align=center]{movem\\la peça verda} (\a+1.9,\a/2);
  \begin{scope}[xshift=4.8cm]
  \fill[fmA!28] (0,0) rectangle (\a,\a-\b);
  \fill[fmK!25] (\a,0) rectangle (\a+\b,\a-\b);
  \draw[fmnavy,line width=1pt] (0,0) rectangle (\a+\b,\a-\b) (\a,0)--(\a,\a-\b);
  \draw[decorate,decoration={brace,amplitude=5pt},thick] (0,\a-\b+0.08)--(\a+\b,\a-\b+0.08) node[midway,above=5pt]{$a+b$};
  \draw[decorate,decoration={brace,amplitude=5pt,mirror},thick] (\a+\b+0.08,0)--(\a+\b+0.08,\a-\b) node[midway,right=5pt]{$a-b$};
  \node[font=\comic,text=fmmuted] at (\a/2+\b/2,-0.4) {$(a+b)(a-b)$};
  \end{scope}
\end{tikzpicture}
\end{minipage}

\subsection{Com s'apliquen}
\begin{exemple}[Exemple: \ primer busquem qui fa de \cA{$a$} i qui fa de \cB{$b$}]
\vspace{-6pt}
\begin{align*}
(\cA{3x}+\cB{2})^2&=(\cA{3x})^2+2\cdot\cA{3x}\cdot\cB{2}+\cB{2}^2=\boxed{9x^2+12x+4} && \text{\nota{$a=3x$,\ $b=2$}}\\
(\cA{x^2}-\cB{5})^2&=(\cA{x^2})^2-2\cdot\cA{x^2}\cdot\cB{5}+\cB{5}^2=\boxed{x^4-10x^2+25}\\
(\cA{2x}+\cB{3y})(\cA{2x}-\cB{3y})&=(\cA{2x})^2-(\cB{3y})^2=\boxed{4x^2-9y^2}
\end{align*}
\end{exemple}

\begin{minipage}[t]{0.5\linewidth}
\begin{compte}[Errors típics]
\mal{(x+3)^2=x^2+9}\ \ {\small falta el doble!}\\[2pt]
\be{(x+3)^2=x^2+6x+9}\\[5pt]
\mal{(3x)^2=3x^2}\ \ {\small el 3 també s'eleva}\\[2pt]
\be{(3x)^2=9x^2}
\end{compte}
\end{minipage}\hfill
\begin{minipage}[t]{0.47\linewidth}
\begin{truc}[També al revés]
Si veus dos quadrats i el doble producte\dots\\[4pt]
$x^2+14x+49=(x+7)^2$\\[2pt]
$x^2-16=(x+4)(x-4)$
\end{truc}
\end{minipage}

% ============================================================ 4.8
\clearpage
\seccio{sec8}{Extracció de factor comú}{Traiem fora el que es repeteix en tots els termes.}

\begin{idea}
\centering{\Large$\cR{a}\cdot b+\cR{a}\cdot c=\cR{a}\,(b+c)$}\\[3pt]
\raggedright Convertim una \textbf{suma} en un \textbf{producte}: diem que \textbf{factoritzem}.
\end{idea}

\begin{minipage}[c]{0.45\linewidth}
\pas{1} \textbf{Nombres}: el \fmkey{MCD} dels coeficients.\\[6pt]
\pas{2} \textbf{Lletres}: les que estan en \textbf{tots} els termes, amb l'exponent \fmkey{més xicotet}.\\[6pt]
\pas{3} \textbf{Dins del parèntesi}: cada terme dividit entre el factor comú.
\end{minipage}\hfill
\begin{minipage}[c]{0.52\linewidth}\centering
\begin{tikzpicture}[t/.style={font=\Large,inner sep=3pt}]
  \node[t] (a) at (0,0) {$12x^3$}; \node[t] (b) at (2.4,0) {$-18x^2$}; \node[t] (c) at (4.8,0) {$+6x$};
  \node[t] (a2) at (0,-1.3) {$\colorbox{fmyellow!70}{$6x$}\cdot2x^2$};
  \node[t] (b2) at (2.4,-1.3) {$\colorbox{fmyellow!70}{$6x$}\cdot(-3x)$};
  \node[t] (c2) at (4.8,-1.3) {$\colorbox{fmyellow!70}{$6x$}\cdot1$};
  \foreach \i in {a,b,c} \draw[-{Latex[length=2.2mm]},fmmuted,line width=1pt] (\i) -- (\i2);
  \node[t,font=\LARGE] (r) at (2.4,-2.75) {$\colorbox{fmyellow!70}{$6x$}\,(2x^2-3x+1)$};
  \node[text=fmgold,font=\fontsize{30}{30}\selectfont] at (-0.6,-2.65) {\faUmbrella};
\end{tikzpicture}
\end{minipage}

\begin{minipage}[t]{0.6\linewidth}
\begin{exemple}
\vspace{-6pt}
\begin{align*}
5x^2-15x&=\boxed{5x\,(x-3)}\\
8x^4+4x^3-12x^2&=\boxed{4x^2\,(2x^2+x-3)}\\
2x\cR{(x-5)}+3\cR{(x-5)}&=\boxed{(x-5)(2x+3)} && \text{\nota{el comú és $(x-5)$}}
\end{align*}
\end{exemple}
\end{minipage}\hfill
\begin{minipage}[t]{0.37\linewidth}
\begin{truc}[Comprova-ho]
Multiplica i has de tornar a l'inici:\\[4pt]
$6x\,(2x^2-3x+1)$\\
$=12x^3-18x^2+6x$ \ \be{}
\end{truc}
\end{minipage}

\subsection{Factoritzar al màxim}
{\large Primer traiem \textbf{factor comú} i després mirem si el que queda és una \textbf{identitat notable}.}

\begin{center}
\begin{tikzpicture}[t/.style={font=\Large,fill=fmdefbg,draw=fmblue,rounded corners=2mm,inner sep=6pt},l/.style={font=\comic\small,text=fmmuted,above,align=center}]
  \node[t] (a) at (0,0) {$x^3-9x$};
  \node[t] (b) at (4.6,0) {$x\,(x^2-9)$};
  \node[t,fill=fmK!14,draw=fmK] (c) at (9.6,0) {$x\,(x+3)(x-3)$};
  \draw[-{Latex[length=3mm]},line width=1.4pt,fmmuted] (a) -- node[l]{factor\\comú} (b);
  \draw[-{Latex[length=3mm]},line width=1.4pt,fmmuted] (b) -- node[l]{suma per\\diferència} (c);
  \node[t] (a) at (0,-2.1) {$2x^2+12x+18$};
  \node[t] (b) at (4.6,-2.1) {$2\,(x^2+6x+9)$};
  \node[t,fill=fmK!14,draw=fmK] (c) at (9.6,-2.1) {$2\,(x+3)^2$};
  \draw[-{Latex[length=3mm]},line width=1.4pt,fmmuted] (a) -- node[l]{factor\\comú} (b);
  \draw[-{Latex[length=3mm]},line width=1.4pt,fmmuted] (b) -- node[l]{quadrat\\d'una suma} (c);
\end{tikzpicture}
\end{center}

% ============================================================ 4.9
\clearpage
\seccio{sec9}{Divisió de polinomis}{Es fa com la divisió de nombres\dots{} però amb $x$.}

\begin{minipage}[c]{0.34\linewidth}\centering
{\Large$\begin{array}{r|l}\cA{17} & \cB{5}\\ \cline{2-2} \cR{2} & \cK{3}\end{array}$}\\[6pt]
{\Large$\cA{17}=\cB{5}\cdot\cK{3}+\cR{2}$}
\end{minipage}\hfill
\begin{minipage}[c]{0.05\linewidth}\centering{\Huge\color{fmmuted}\faArrowRight}\end{minipage}\hfill
\begin{minipage}[c]{0.57\linewidth}
\begin{idea}
\centering{\Large$\cA{D(x)}=\cB{d(x)}\cdot\cK{C(x)}+\cR{R(x)}$}\\[4pt]
\raggedright\small\cA{\textbf{dividend}} \ = \ \cB{\textbf{divisor}} $\cdot$ \cK{\textbf{quocient}} + \cR{\textbf{residu}}\\[2pt]
Acabem quan el grau del \cR{residu} és \textbf{menor} que el del \cB{divisor}. Si el residu és 0, la divisió és \textbf{exacta}.
\end{idea}
\end{minipage}

\subsection{Pas a pas}
\begin{minipage}[c]{0.4\linewidth}
\pas{1} \textbf{Ordenem} els dos polinomis.\\[5pt]
\pas{2} \textbf{Dividim} el primer terme del dividend entre el primer del divisor.\\[5pt]
\pas{3} \textbf{Multipliquem} per tot el divisor i ho escrivim davall \textbf{canviat de signe}. \textbf{Sumem}.\\[5pt]
\pas{4} \textbf{Repetim} fins que no es puga més.
\end{minipage}\hfill
\begin{minipage}[c]{0.57\linewidth}
\begin{exemple}[Exemple: \ $(x^3-2x^2+4x-5):(x-2)$]
\renewcommand{\arraystretch}{1.35}\large
\[\begin{array}{rrrr|l@{\qquad}l}
x^3 & -2x^2 & +4x & -5 & \,\cB{x-2} & \\ \cline{5-5}
\cR{-x^3} & \cR{+2x^2} & & & \,\cK{x^2+4} & \text{\nota{\pas{2}\ $x^3:x=\cK{x^2}$}}\\ \cline{1-2}
 & 0 & +4x & -5 & & \text{\nota{\pas{3}\ sumem}}\\
 & & \cR{-4x} & \cR{+8} & & \text{\nota{\pas{4}\ $4x:x=\cK{+4}$}}\\ \cline{3-4}
 & & & \cR{\boxed{3}} & & \text{\nota{residu}}
\end{array}\]
\normalsize $\cK{C(x)=x^2+4}$\qquad $\cR{R=3}$
\end{exemple}
\end{minipage}

\begin{minipage}[t]{0.52\linewidth}
\begin{compte}[Falta algun grau?]
Deixa un \textbf{buit} (o posa-hi un 0):\\[3pt]
$x^4\qquad\ \ -3x^2+2x$\ \ \ \notab{falten $x^3$ i el nombre}
\end{compte}
\end{minipage}\hfill
\begin{minipage}[t]{0.45\linewidth}
\begin{truc}[Comprova-ho]
$\cB{(x-2)}\cdot\cK{(x^2+4)}+\cR{3}$\\
$=x^3-2x^2+4x-8+3$\\
$=x^3-2x^2+4x-5$ \ \be{}
\end{truc}
\end{minipage}

% ============================================================ 4.10
\clearpage
\seccio{sec10}{Regla de Ruffini}{Una drecera per a dividir entre \,$(x-a)$.}

\begin{idea}
Quan el divisor és \fmkey{$x-a$}, podem dividir \textbf{molt més ràpid}: treballem \textbf{només amb els coeficients}.
\end{idea}

{\large Exemple: \ $(2x^3-5x^2+4):(x-3)$}

\begin{minipage}[c]{0.44\linewidth}\centering
\ruffini{3}{2,-5,0,4}{6,3,9}{2,1,3,13}{4}\\[4pt]
\tcbox[colback=fmK!10,colframe=fmK,arc=3mm,boxrule=1.2pt,left=6pt,right=6pt]{\large$C(x)=2x^2+x+3$\qquad\cR{$R=13$}}\par
{\comic\small\color{fmmuted}el quocient té un grau menys}
\end{minipage}\hfill
\begin{minipage}[c]{0.53\linewidth}
\pas{1} Escrivim els \textbf{coeficients} (amb \textbf{0} on falta un terme). A l'esquerra, el nombre \cB{$a$}.\\[5pt]
\pas{2} {\color{fmK}\faArrowDown}\ \textbf{Baixem} el primer coeficient.\\[5pt]
\pas{3} {\color{fmB}\faArrowRight}\ \textbf{Multipliquem} per \cB{$a$} i ho posem a la columna següent.\\[5pt]
\pas{4} {\color{fmA}\faArrowDown}\ \textbf{Sumem}. I tornem a \pas{3}.\\[5pt]
\pas{5} L'\cR{\textbf{últim}} és el \cR{\textbf{residu}}; els altres, el quocient.
\end{minipage}

\begin{compte}[Compte amb el signe de $a$!]
\begin{tabular}{@{}l@{\qquad}l@{}}
$x\cB{-3}$ \ $\to$ \ posem \ \cB{\textbf{3}} & $x\cB{+2}=x-(-2)$ \ $\to$ \ posem \ \cB{$\mathbf{-2}$}
\end{tabular}
\end{compte}

\subsection{Teorema del residu}
\begin{minipage}[c]{0.5\linewidth}
\begin{idea}
\centering{\Large El residu de \ $P(x):(x-a)$ \ és \ $\boxed{R=P(a)}$}\\[4pt]
\raggedright\small $P(3)=2\cdot27-5\cdot9+4=13$ \ \be{} (és el mateix residu!)
\end{idea}
\end{minipage}\hfill
\begin{minipage}[c]{0.47\linewidth}
\begin{truc}[Si el residu és 0\dots]
$P(a)=0$ \ $\Rightarrow$ \ la divisió és \textbf{exacta}:
\begin{itemize}[leftmargin=1.1em,itemsep=0pt]
\item $a$ és una \textbf{arrel} de $P$,
\item $(x-a)$ és un \textbf{factor}: \ $P(x)=(x-a)\cdot C(x)$.
\end{itemize}
\end{truc}
\end{minipage}

\begin{exemple}[Exemple: \ $(x^3+2x^2-5x-6):(x-2)$]
\begin{minipage}[c]{0.42\linewidth}\centering
\ruffini{2}{1,2,-5,-6}{2,8,6}{1,4,3,0}{4}
\end{minipage}\hfill
\begin{minipage}[c]{0.55\linewidth}
$C(x)=x^2+4x+3$,\quad \cR{$R=0$}\ \ \nota{exacta!}\\[6pt]
{\large$x^3+2x^2-5x-6=(x-2)(x^2+4x+3)$}\\[4pt]
\nota{2 és una arrel i $(x-2)$ n'és un factor}
\end{minipage}
\end{exemple}

% ============================================================ resum
\clearpage
{\color{fmnavy}\fontsize{23}{27}\heavy\selectfont\faClipboardList\ \ En resum}\par\vspace{5pt}{\color{fmblue}\titlerule[2pt]}
\frase{Tota la unitat en una pàgina.}
\newcommand{\rtile}[3]{\begin{tcolorbox}[enhanced,colback=white,colframe=fmblue!45,arc=3mm,boxrule=1pt,left=3mm,right=3mm,top=2.5mm,bottom=2.5mm,fontupper=\large,
  title={\tcbox[on line,colback=fmblue,colframe=fmblue,arc=2mm,boxsep=0pt,left=4pt,right=4pt,top=1pt,bottom=1pt]{\color{white}\small\heavy #1}\ \ #2},coltitle=fmnavy,colbacktitle=fmdefbg,fonttitle=\semi,toptitle=0.5mm,bottomtitle=0.5mm]#3\end{tcolorbox}}
\begin{tcbraster}[raster columns=2,raster equal height=rows,raster column skip=14pt,raster row skip=22pt,raster before skip=10pt]
\rtile{4.1}{Llenguatge algebraic}{El doble: $2x$. \ El següent: $x+1$. \\
\ Un senar: $2x+1$.\ Un parell $2x$\\ [2ex]
\textbf{Valor numèric}: canvia la lletra pel nombre \textbf{(entre parèntesis)}.}
\rtile{4.2}{Monomis}{$\cA{-5}\,\cB{x^2y}$: \ \cA{coeficient} $-5$, \ \cB{part literal} $x^2y$, \ \cK{grau} 3.\\[2ex] \textbf{Semblants}: la mateixa part literal.}
\rtile{4.3}{Operacions amb monomis}{\textbf{Suma}: només semblants; sumem coeficients: $5x^2+3x^2=8x^2$.\\ \textbf{Producte}: sumem exponents.\\ \textbf{Quocient}: els restem.}
\rtile{4.4}{Polinomis}{Grau = el més alt. Coeficient principal. Terme independent.\\[2ex] $a$ és una \textbf{arrel} si $P(a)=0$.}
\rtile{4.5}{Suma i resta}{En columnes, cada família amb la seua.\\ \textbf{Restar}: canvia el signe de \textbf{tots} els termes del segon.}
\rtile{4.6}{Producte}{Tots per tots i reduïm.\\ Grau del producte = suma dels graus.}
\rtile{4.7}{Identitats notables}{$(a\pm b)^2=a^2\pm2ab+b^2$\\[2ex]$(a+b)(a-b)=a^2-b^2$}
\rtile{4.8}{Factor comú}{MCD dels nombres i lletres comunes amb l'exponent més xicotet:\\ $12x^3-18x^2+6x=6x\,(2x^2-3x+1)$}
\rtile{4.9}{Divisió}{$D=d\cdot C+R$, \ amb grau $R<$ grau $d$.\\ Si $R=0$, la divisió és exacta.}
\rtile{4.10}{Ruffini i teorema del residu}{Per a dividir entre $x-a$: baixa, multiplica, suma.\\ El residu és $R=P(a)$.}
\end{tcbraster}
\end{document}
